\subsection{\texorpdfstring{Relations between \(\mathbf{S}(\mathbf{M}^{*})\), \(\mathbf{TS}(\mathbf{M})^{*\mathrm{gr}}\) and \(\operatorname{Pol}(\mathbf{M}, \mathbf{A})\)}{Relations between S(M*), the graded dual of TS(M), and Pol(M, A)}}

Let M be a free A-module, we shall equip the graded dual \(\mathbf{TS}(\mathbf{M})^{*gr}\) with the structure of a commutative, associative and unital graded \(^{1}\) algebra, derived from the graded cogebra structure of \(\mathbf{TS}(\mathbf{M})\) (III, p.~580). By III, p.~497 there exists a unique homomorphism of graded A-algebras

\[
\theta : \mathbf {S} (\mathbf {M} ^ {*}) \rightarrow \mathbf {T S} (\mathbf {M}) ^ {* \mathrm{gr}}
\]

inducing in degree 1 the identity mapping of \(\mathbf{M}^*\) .

PROPOSITION 20. --- If the A-module M is free and finitely generated, then \(\theta\) is an isomorphism of graded algebras.

Let \((e_i)_{i\in I}\) be a basis of \(\mathbf{M}\) and \((e_i^*)_{i\in I}\) the dual basis of \(\mathbf{M}^*\) . For \(\nu \in \mathbf{N}^I\) put

\[
e _ {\nu} = \prod_ {i \in I} \gamma_ {\nu_ {i}} (e _ {i}) \in \mathbf {T S} (\mathbf {M}).
\]

By Prop. 4 (IV, p.~47) the family \((e_{\nu})_{\nu \in \mathbf{N}^{I}}\) is a basis of \(\mathbf{TS}(\mathbf{M})\) ; let \((e_{\nu}^{*})\) be the basis of \(\mathbf{TS}(\mathbf{M})^{*\mathrm{gr}}\) dual to \((e_{\nu})\) . In the light of III, p.~505, Th. 1 it is enough to show that for any \(\nu \in \mathbf{N}^{I}\) we have

\[
e _ {\nu} ^ {*} = \prod_ {i \in I} \left(e _ {i} ^ {*}\right) ^ {\nu_ {i}},
\]

or also that for \(\rho, \sigma \in \mathbf{N}^{\mathrm{I}}\) we have \(e_{\rho}^{*} \cdot e_{\sigma}^{*} = e_{\rho + \sigma}^{*}\) ; but this last assertion follows from IV, p.~50, Formula (12).

Remark 1. --- In the same way we see that if M is a finitely generated free A-module, then the graded algebra \(\mathbf{S}(\mathbf{M})^{*\mathrm{gr}}\) defined in III, p.~593, may be identified with \(\mathbf{TS}(\mathbf{M}^*)\) .

PROPOSITION 21. --- The canonical homomorphism of A-modules (IV, p.~55)

\[
u: \mathbf {T S} (\mathrm{M}) ^ {* \mathrm{gr}} \rightarrow \operatorname{Pol} _ {\mathrm{A}} (\mathrm{M}, \mathrm{A})
\]

is an algebra homomorphism.

Let \(a \in TS^{q}(M)^{*}\) , \(b \in TS^{r}(M)^{*}\) , \(x \in M\) ; we have

\[
\begin{array}{r l} u (a b) (x) & = \langle a b, \gamma_ {q + r} (x) \rangle = \langle a \otimes b, c (\gamma_ {q + r} (x)) \rangle = \\ & = \langle a \otimes b, \gamma_ {q} (x) \otimes \gamma_ {r} (x) \rangle = \langle a, \gamma_ {q} (x) \rangle \langle b, \gamma_ {r} (x) \rangle = u (a) (x). u (b) (x), \end{array}
\]

whence the result.

Remarks. --- 2) The composite homomorphism \(\lambda_{\mathbf{M}}=u\circ\theta:\mathbf{S}(\mathbf{M}^{*})\to\mathrm{Pol}_{\mathbf{A}}(\mathbf{M},\mathbf{A})\) is the unique unital homomorphism of algebras inducing the inclusion of

\[
\mathbf {M} ^ {*} = \operatorname{Pol} ^ {1} (\mathbf {M}, \mathbf {A})
\]

in \(\operatorname{Pol}(\mathbf{M}, \mathbf{A})\) . If \(\mathbf{M}\) is finitely generated free and \(\mathbf{A}\) is an infinite integral domain, then \(\lambda_{\mathbf{M}}\) is bijective (Prop. 20 and Cor. of Prop. 19). In particular if \(\mathbf{A}\) is an infinite integral domain, then the canonical homomorphism \(f \mapsto \tilde{f}\) of \(\mathbf{A}[\mathbf{X}_1, \ldots, \mathbf{X}_n]\) into \(\operatorname{Pol}(\mathbf{A}^n, \mathbf{A})\) (IV, p.~4) is an isomorphism.

\begin{enumerate}
\def\labelenumi{\arabic{enumi})}
\setcounter{enumi}{2}
\tightlist
\item
  Consider the coproduct \(c_{\mathbf{S}}: \mathbf{S}(\mathbf{M}^{*}) \to \mathbf{S}(\mathbf{M}^{*} \times \mathbf{M}^{*})\) (III, p.~575, Example 6). For any \(v \in \mathbf{S}(\mathbf{M}^{*})\) , \(x, y \in \mathbf{M}\) , the polynomial mapping \(\lambda_{\mathbf{M} \times \mathbf{M}}(c_{\mathbf{S}}(v)): \mathbf{M} \times \mathbf{M} \to \mathbf{A}\) maps \((x, y)\) to \(\lambda_{\mathbf{M}}(v)(x + y)\) . For the two algebra homomorphisms of \(\mathbf{S}(\mathbf{M}^{*})\) into \(\operatorname{Map}(\mathbf{M} \times \mathbf{M}, \mathbf{A})\) defined in this way agree on \(\mathbf{M}^{*}\) , in virtue of the relation
\end{enumerate}

\[
(v \otimes 1 + 1 \otimes v) (x, y) = v (x + y) \quad (v \in \mathbf {M} ^ {*}).
\]

